\chapter{Charakterisierungen und Randfälle}
\label{b08guide:ch:criteria}

\section{Eindeutige Umkehrbarkeit}

\begin{definition}[Umgedrehter Graph und beidseitige Umkehrfunktion]
\label{b08guide:def:inverse}
Für $f\colon A\to B$ heißt
\[
 f^{\leftarrow}=\{(b,a)\in B\times A\mid f(a)=b\}
\]
der \emph{umgedrehte Graph}. Eine Funktion $g\colon B\to A$
heißt \emph{beidseitige Umkehrfunktion} von $f$, wenn
\[
 g\circ f=\operatorname{id}_A,
 \qquad f\circ g=\operatorname{id}_B.
\]
\end{definition}

\begin{theorem}[Charakterisierung durch Umkehrbarkeit]
\label{b08guide:thm:inversecriteria}
Für eine Funktion $f\colon A\to B$ sind äquivalent:
\begin{enumerate}
\item $f$ ist bijektiv von $A$ auf $B$.
\item Für jedes $b\in B$ besteht $f^{-1}[\{b\}]$ aus genau einem Element.
\item Der umgedrehte Graph $f^{\leftarrow}$ ist eine Funktion von $B$ nach $A$.
\item $f$ besitzt eine beidseitige Umkehrfunktion.
\item Es gibt Funktionen $r,s\colon B\to A$ mit
      $r\circ f=\operatorname{id}_A$ und $f\circ s=\operatorname{id}_B$.
\end{enumerate}
In diesem Fall sind die beidseitige Umkehrfunktion und auch die
Funktionen $r,s$ aus Bedingung~5 eindeutig bestimmt und gleich
$f^{\leftarrow}$. Diese Funktion wird mit $f^{-1}$ bezeichnet.
Für $a\in A$ und $b\in B$ gilt
\[
 f(a)=b\quad\Longleftrightarrow\quad f^{-1}(b)=a.
\]
Die Aussage gilt auch für leere Mengen und benötigt kein Auswahlaxiom.
\end{theorem}

\section{Erkennen an sämtlichen Teilmengen}

\begin{theorem}[Potenzmengen- und Rückgewinnungskriterien]
\label{b08guide:thm:powersetcriteria}
Für $f\colon A\to B$ sind äquivalent:
\begin{enumerate}
\item $f\colon A\bij B$.
\item $f_*\colon\mathcal P(A)\to\mathcal P(B)$ ist bijektiv.
\item $f^*\colon\mathcal P(B)\to\mathcal P(A)$ ist bijektiv.
\item Beide Rückgewinnungsregeln gelten:
\[
 \begin{aligned}
 f^{-1}[f[U]]&=U &&(U\subseteq A),\\
 f[f^{-1}[V]]&=V &&(V\subseteq B).
 \end{aligned}
\]
\end{enumerate}
In Bedingung~4 genügt die erste Regel für alle Einermengen
$U=\{a\}$ zusammen mit der zweiten Regel für $V=B$.
Die Kriterien~2 und~3 beziehen sich auf die von $f$ induzierten
Abbildungen, nicht auf eine beliebige Bijektion zwischen Potenzmengen.
\end{theorem}

\begin{theorem}[Komplementkriterium mit der ganzen Zielmenge]
\label{b08guide:thm:complementcriteria}
Für $f\colon A\to B$ gilt
\[
 f\colon A\bij B
 \quad\Longleftrightarrow\quad
 \forall U\subseteq A\quad f[A\setminus U]=B\setminus f[U].
\]
Dabei muß auch $U=\varnothing$ zugelassen werden.
Auf der rechten Seite wird das Komplement ausdrücklich in $B$ gebildet.
\end{theorem}

\section{Leere Mengen und Einermengen}

\begin{theorem}[Bijektionen in den kleinsten Fällen]
\label{b08guide:thm:empty}
Für jede Menge $A$ gilt
\[
 \operatorname{Bij}(\varnothing,A)
 =\operatorname{Bij}(A,\varnothing)
 =\begin{cases}
   \{\varnothing\},&A=\varnothing,\\
   \varnothing,&A\ne\varnothing.
  \end{cases}
\]
Die einzige Bijektion zwischen leeren Mengen ist also der leere Graph;
er ist seine eigene Umkehrfunktion. Für jedes Element $b$ gilt ferner
\[
 \operatorname{Bij}(A,\{b\})
 =\begin{cases}
   \{\{(a,b)\}\},& A=\{a\}\text{ für ein }a,\\
   \varnothing,& A\text{ ist keine Einermenge}.
  \end{cases}
\]
Insbesondere ist $a\mapsto b$ die einzige Bijektion
$\{a\}\bij\{b\}$.
\end{theorem}
